
% ============================================================
% Appendix: Theoretical Guarantees for Contextual Hamiltonian Enrichment
% Drop this file into the NeurIPS source and include with:
%   \input{appendix_theory}
%
% Required packages in the main preamble:
%   \usepackage{amsmath,amssymb,amsthm,mathtools}
%   \usepackage{natbib}
%
% Optional theorem environments if not already defined:
%   \newtheorem{theorem}{Theorem}
%   \newtheorem{lemma}{Lemma}
%   \newtheorem{proposition}{Proposition}
%   \newtheorem{corollary}{Corollary}
%   \theoremstyle{definition}
%   \newtheorem{assumption}{Assumption}
%   \newtheorem{remark}{Remark}
% ============================================================

\section{Theoretical Guarantees for Contextual Hamiltonian Enrichment}
\label{app:theory}

This appendix collects the formal properties used in Section~3.6.
The goal is not to claim unconditional safety or global optimality. Instead, we prove
that contextual Hamiltonian enrichment has five conditional properties aligned with
the central claims of the paper: (i) an added energy term induces a new force and
sensitivity channel; (ii) the enriched system preserves the port-Hamiltonian
dissipativity structure; (iii) the discrete rollout is practically dissipative for small
enough step size; (iv) when context gradients vanish, the context-enriched field remains close to the
geometry scaffold and cannot generate a spurious lateral escape; and (v) when a
projected risk-gradient margin exists, the contextual force activates and locally
reduces risk. We also justify the detached-quantile CVaR gradient used in the
episode loop. The passivity arguments follow the standard port-Hamiltonian energy
balance~\citep{vanDerSchaftJeltsema2014,Desai2021PortHNN}; the CVaR argument
uses the Rockafellar--Uryasev variational representation~\citep{RockafellarUryasev2000}
and standard empirical quantile sensitivity arguments~\citep{HongHuLiu2014}.





\subsection{Setup and notation}
\label{app:setup}

Let \(q \in \mathbb R^d\) denote the configuration coordinate and \(p\in\mathbb R^d\)
its cotangent variable. Let \(x=(q,p)\), let \(M\succ 0\) be the mass matrix, and
write
\[
    v = M^{-1}p .
\]
The geometry-scaffold Hamiltonian is
\[
    H_1(q,p)
    =
    \frac12 p^\top M^{-1}p + H_{\mathrm{geom}}(q),
\]
where \(H_{\mathrm{geom}}\) contains the goal and geometric barrier terms. The context-enriched field
adds the context Hamiltonian
\[
    H_{\mathrm{ctx}}(q;z)
    =
    \lambda_s \widetilde r(\mathbf c;z)
    +
    \lambda_h b_{\mathrm{sp}}(\phi(\mathbf c;z)),
\]
where \(\mathbf c\) is the workspace position component of \(q\), \(z\) denotes the
locally sensed context patch, \(\widetilde r\) is the smoothed soft-risk field, and
\(\phi\) is the signed-distance field for hard hazards. Thus
\[
    H_2(q,p;z)
    =
    H_1(q,p)+H_{\mathrm{ctx}}(q;z).
\]
The corresponding context force is
\[
    F_{\mathrm{ctx}}(q;z)
    =
    -\nabla_q H_{\mathrm{ctx}}(q;z)
    =
    -\lambda_s \nabla_q \widetilde r(\mathbf c;z)
    -
    \lambda_h b_{\mathrm{sp}}'(\phi(\mathbf c;z))\nabla_q \phi(\mathbf c;z).
\]
This is exactly the soft-risk plus hard-hazard force decomposition used in the main
paper.

The continuous-time enriched port-Hamiltonian dynamics are
\begin{align}
    \dot q &= \nabla_p H_2(q,p;z)=M^{-1}p, \label{eq:ct_q}\\
    \dot p &=
    -\nabla_q H_2(q,p;z)
    -
    D(q,p;z)M^{-1}p
    +
    B(q,p;z)u , \label{eq:ct_p}
\end{align}
where \(D(q,p;z)\succeq 0\) is the dissipative port. The rollout used in the paper is
the semi-implicit update
\begin{align}
    p_{t+1}
    &=
    p_t
    -\tau\nabla_q H_{\mathrm{geom}}(q_t)
    -\tau\nabla_q H_{\mathrm{ctx}}(q_t;z_t)
    -\tau D_tM^{-1}p_t
    +\tau B_tu_t, \label{eq:disc_p}\\
    q_{t+1}
    &=
    q_t+\tau M^{-1}p_{t+1}. \label{eq:disc_q}
\end{align}
When \(D_tM^{-1}p_t\) is implemented as the scalar dissipative port
\(R_\theta(q_t,p_t)=\gamma M^{-1}p_t\), this reduces to the update in the main text.

\begin{assumption}[Rollout regularity]
\label{ass:regularity}
There exists a compact rollout tube \(\mathcal K\subset \mathbb R^{2d}\) such that
the geometry scaffold and context-enriched field trajectories remain in \(\mathcal K\) for the horizon under
analysis. On \(\mathcal K\), \(H_{\mathrm{geom}}\) and \(H_{\mathrm{ctx}}\) are \(C^2\),
and their gradients are Lipschitz.
\end{assumption}

\begin{assumption}[Bounded context force]
\label{ass:bounded_ctx}
There exists \(G_{\mathrm{ctx}}<\infty\) such that
\[
    \|\nabla_q H_{\mathrm{ctx}}(q;z)\|\le G_{\mathrm{ctx}}
\]
for all \((q,p)\in\mathcal K\) and all context patches considered.
\end{assumption}

\begin{assumption}[Dissipative port]
\label{ass:diss}
There exists \(\rho\ge 0\) such that
\[
    D(q,p;z)\succeq \rho I
\]
on \(\mathcal K\). When \(\rho>0\), the system is strictly dissipative in velocity.
\end{assumption}

\begin{assumption}[Local Lipschitz rollout map]
\label{ass:lipschitz_rollout}
The shared geometry-scaffold update map is locally Lipschitz on \(\mathcal K\): there exists
\(L_\Phi\ge 0\) such that one step of the shared dynamics amplifies state errors by at
most \(1+\tau L_\Phi\).
\end{assumption}

\subsection{Energy enrichment induces force and sensitivity channels}

\begin{lemma}[Energy enrichment induces force, sensitivity, and excitation channels]
\label{lem:enrichment_channels}
Let
\[
    H_2(q,p;\theta,\vartheta)
    =
    H_1(q,p;\theta)+\vartheta H_{\mathrm{new}}(q),
\]
where \(H_{\mathrm{new}}\in C^1\) and \(\vartheta\in\mathbb R\) is a learnable scalar.
Then:
\begin{enumerate}
    \item the enriched momentum equation contains the new force channel
    \[
        F_{\mathrm{new}}(q)=-\vartheta\nabla_q H_{\mathrm{new}}(q);
    \]
    \item differentiating a rollout-dependent loss \(J\) through the discrete update
    creates the direct sensitivity stream
    \[
        \frac{\partial J}{\partial \vartheta}
        =
        -\tau\sum_t
        \left\langle
        \frac{\partial J}{\partial p_{t+1}},
        \nabla_qH_{\mathrm{new}}(q_t)
        \right\rangle
        +
        \textnormal{recursive terms};
    \]
    \item local identifiability of \(\vartheta\) from force observations requires
    non-degenerate excitation of the new channel, i.e. the empirical Gram matrix
    over \(\{\nabla_qH_{\mathrm{new}}(q_t)\}_t\) must be nonsingular on the
    observed force subspace.
\end{enumerate}
\end{lemma}

\begin{proof}
Since
\[
    \nabla_q H_2
    =
    \nabla_q H_1+\vartheta\nabla_qH_{\mathrm{new}},
\]
Hamilton's momentum equation contains
\[
    -\nabla_qH_2
    =
    -\nabla_qH_1-\vartheta\nabla_qH_{\mathrm{new}}.
\]
Thus the added term contributes the force channel
\(F_{\mathrm{new}}=-\vartheta\nabla_qH_{\mathrm{new}}\).

For the sensitivity statement, consider the discrete momentum update
\[
    p_{t+1}
    =
    p_t-\tau\nabla_qH_1(q_t,p_t)
    -\tau\vartheta\nabla_qH_{\mathrm{new}}(q_t)+\cdots .
\]
Differentiating with respect to \(\vartheta\) gives
\[
    \frac{\partial p_{t+1}}{\partial\vartheta}
    =
    \frac{\partial p_t}{\partial\vartheta}
    -\tau\nabla_qH_{\mathrm{new}}(q_t)
    -\tau\nabla_q^2H_2(q_t,p_t)
        \frac{\partial q_t}{\partial\vartheta}
    +\cdots .
\]
Applying the chain rule to any rollout loss \(J\) yields
\[
    \frac{\partial J}{\partial\vartheta}
    =
    \sum_t
    \left\langle
    \frac{\partial J}{\partial p_{t+1}},
    \frac{\partial p_{t+1}}{\partial\vartheta}
    \right\rangle
    +
    \sum_t
    \left\langle
    \frac{\partial J}{\partial q_{t+1}},
    \frac{\partial q_{t+1}}{\partial\vartheta}
    \right\rangle .
\]
The displayed direct stream is the term produced by
\(-\tau\nabla_qH_{\mathrm{new}}(q_t)\); the remaining terms are recursive
contributions through \(\partial q_t/\partial\vartheta\) and
\(\partial p_t/\partial\vartheta\).

For identifiability, suppose two coefficients \(\vartheta,\vartheta'\) induce the same
observed force on the rollout. Then
\[
    (\vartheta-\vartheta')\nabla_qH_{\mathrm{new}}(q_t)=0
\]
in all observed force directions. A nonzero coefficient difference is impossible only
when the observed gradients span the relevant force subspace. Equivalently, the
corresponding empirical Gram matrix is nonsingular on that subspace.
\end{proof}

\subsection{Passivity and practical discrete dissipativity}

\begin{theorem}[Contextual Hamiltonian enrichment preserves dissipativity]
\label{thm:continuous_passivity}
Assume~\ref{ass:regularity}--\ref{ass:diss}. If the context \(z\) is
fixed and \(u=0\), then along solutions of \eqref{eq:ct_q}--\eqref{eq:ct_p},
\[
    \frac{d}{dt}H_2(q(t),p(t);z)
    =
    -v(t)^\top D(q(t),p(t);z)v(t)
    \le
    -\rho\|v(t)\|^2
    \le 0 .
\]
Thus the addition of \(H_{\mathrm{ctx}}\) preserves the port-Hamiltonian dissipativity
structure. If \(z=z(t)\) and \(u\neq0\), then
\[
    \frac{d}{dt}H_2
    =
    -v^\top Dv
    +
    \left\langle \nabla_zH_{\mathrm{ctx}}(q;z),\dot z \right\rangle
    +
    \left\langle \nabla_pH_2,Bu \right\rangle .
\]
Hence the enriched system is passive with storage \(H_2\) and supply rate
\[
    s(t)
    =
    \left\langle \nabla_zH_{\mathrm{ctx}}(q;z),\dot z \right\rangle
    +
    \left\langle \nabla_pH_2,Bu \right\rangle .
\]
\end{theorem}

\begin{proof}
For fixed \(z\), differentiate \(H_2\) along trajectories:
\[
    \dot H_2
    =
    \langle\nabla_qH_2,\dot q\rangle
    +
    \langle\nabla_pH_2,\dot p\rangle .
\]
Using \(\dot q=\nabla_pH_2\) and
\(\dot p=-\nabla_qH_2-D\nabla_pH_2\), we obtain
\[
    \dot H_2
    =
    \langle\nabla_qH_2,\nabla_pH_2\rangle
    +
    \left\langle
    \nabla_pH_2,
    -\nabla_qH_2-D\nabla_pH_2
    \right\rangle .
\]
The two interconnection terms cancel:
\[
    \langle\nabla_qH_2,\nabla_pH_2\rangle
    -
    \langle\nabla_pH_2,\nabla_qH_2\rangle=0.
\]
Therefore
\[
    \dot H_2
    =
    -\langle\nabla_pH_2,D\nabla_pH_2\rangle.
\]
Since \(\nabla_pH_2=M^{-1}p=v\),
\[
    \dot H_2=-v^\top Dv.
\]
By Assumption~\ref{ass:diss}, \(D\succeq \rho I\), so
\[
    -v^\top Dv\le -\rho\|v\|^2\le 0.
\]
If \(z=z(t)\), the chain rule adds
\[
    \langle\nabla_zH_2,\dot z\rangle
    =
    \langle\nabla_zH_{\mathrm{ctx}},\dot z\rangle .
\]
If \(u\neq 0\), the input term \(Bu\) in \(\dot p\) contributes
\[
    \langle\nabla_pH_2,Bu\rangle .
\]
Combining these terms gives the claimed supply-rate identity.
\end{proof}

\begin{theorem}[Practical discrete dissipativity and step-size condition]
\label{thm:discrete_passivity}
Assume~\ref{ass:regularity}--\ref{ass:diss} and \(u_t=0\).
For the semi-implicit context-enriched update \eqref{eq:disc_p}--\eqref{eq:disc_q}, there
exists a finite constant \(C_{\mathcal K}>0\) such that
\[
    H_2(x_{t+1};z_t)-H_2(x_t;z_t)
    \le
    -\tau v_t^\top D_tv_t+C_{\mathcal K}\tau^2 .
\]
Consequently, the discrete rollout is practically dissipative. If
\[
    \tau < \frac{v_t^\top D_tv_t}{C_{\mathcal K}},
\]
then \(H_2(x_{t+1};z_t)<H_2(x_t;z_t)\). If
\(v_t^\top D_tv_t\ge d_{\min}>0\) uniformly on a region of interest, then any
\[
    \tau < \frac{d_{\min}}{C_{\mathcal K}}
\]
guarantees one-step energy decrease on that region.
\end{theorem}

\begin{proof}
Let \(f_z(x)\) be the continuous-time vector field defined by
\eqref{eq:ct_q}--\eqref{eq:ct_p} with \(u=0\). Since the semi-implicit update is a
first-order consistent discretization of this vector field and all derivatives are bounded
on compact \(\mathcal K\), there exists \(C_I>0\) such that
\[
    x_{t+1}=x_t+\tau f_{z_t}(x_t)+e_t,
    \qquad
    \|e_t\|\le C_I\tau^2 .
\]
Because \(H_2\in C^2(\mathcal K)\), its gradient is Lipschitz. Let \(L_H\) be a
Lipschitz constant for \(\nabla H_2\). The smoothness inequality gives
\[
    H_2(x_{t+1})-H_2(x_t)
    \le
    \langle\nabla H_2(x_t),x_{t+1}-x_t\rangle
    +
    \frac{L_H}{2}\|x_{t+1}-x_t\|^2 .
\]
Substituting \(x_{t+1}-x_t=\tau f_{z_t}(x_t)+e_t\),
\begin{align*}
    H_2(x_{t+1})-H_2(x_t)
    &\le
    \tau\langle\nabla H_2(x_t),f_{z_t}(x_t)\rangle
    +
    \langle\nabla H_2(x_t),e_t\rangle  \\
    &\qquad
    +
    \frac{L_H}{2}\|\tau f_{z_t}(x_t)+e_t\|^2 .
\end{align*}
By Theorem~\ref{thm:continuous_passivity},
\[
    \langle\nabla H_2(x_t),f_{z_t}(x_t)\rangle
    =
    -v_t^\top D_tv_t .
\]
The remaining two terms are bounded by \(C_{\mathcal K}\tau^2\) for some finite
constant \(C_{\mathcal K}\), since \(\mathcal K\) is compact and
\(\|\nabla H_2\|\), \(\|f_z\|\), and \(\|e_t\|/\tau^2\) are bounded there. Hence
\[
    H_2(x_{t+1};z_t)-H_2(x_t;z_t)
    \le
    -\tau v_t^\top D_tv_t+C_{\mathcal K}\tau^2 .
\]
Strict decrease follows whenever
\(C_{\mathcal K}\tau^2<\tau v_t^\top D_tv_t\). Dividing by \(\tau>0\) yields the
step-size condition. The uniform-margin result follows by replacing
\(v_t^\top D_tv_t\) by \(d_{\min}\).
\end{proof}

\subsection{Scaffold preservation and no-hallucination}

\begin{theorem}[Scaffold preservation under small context force]
\label{thm:scaffold_preservation}
Let \(x_t^{(1)}=(q_t^{(1)},p_t^{(1)})\) be the geometry-scaffold rollout and
\(x_t^{(2)}=(q_t^{(2)},p_t^{(2)})\) be the context-enriched rollout, initialized from the same
state. Assume~\ref{ass:regularity}, \ref{ass:bounded_ctx}, and
\ref{ass:lipschitz_rollout}. Suppose that along the rollout tube
\[
    \|\nabla_qH_{\mathrm{ctx}}(q;z)\|\le \varepsilon .
\]
Then after \(N\) steps,
\[
    \|x_N^{(2)}-x_N^{(1)}\|
    \le
    \tau\sqrt{1+\tau^2\|M^{-1}\|^2}\,
    \varepsilon\,
    \frac{(1+\tau L_\Phi)^N-1}{\tau L_\Phi},
\]
with the last factor interpreted as \(N\) when \(L_\Phi=0\). In particular, if
\(\varepsilon=0\), the context-enriched rollout recovers the geometry scaffold up to numerical
roundoff.
\end{theorem}

\begin{proof}
The only difference between the geometry scaffold and context-enriched field is the context term in the
momentum update. At one step, this perturbation is
\[
    \delta p_{t+1}
    =
    -\tau\nabla_qH_{\mathrm{ctx}}(q_t;z_t),
\]
and the induced position perturbation is
\[
    \delta q_{t+1}
    =
    \tau M^{-1}\delta p_{t+1}
    =
    -\tau^2M^{-1}\nabla_qH_{\mathrm{ctx}}(q_t;z_t).
\]
Thus
\[
    \|\delta p_{t+1}\|\le \tau\varepsilon,
    \qquad
    \|\delta q_{t+1}\|\le \tau^2\|M^{-1}\|\varepsilon.
\]
The one-step state perturbation is therefore bounded by
\[
    \|\delta x_{t+1}\|
    \le
    \tau\sqrt{1+\tau^2\|M^{-1}\|^2}\,\varepsilon .
\]
Let
\[
    e_t=\|x_t^{(2)}-x_t^{(1)}\|.
\]
By Assumption~\ref{ass:lipschitz_rollout},
\[
    e_{t+1}
    \le
    (1+\tau L_\Phi)e_t
    +
    \tau\sqrt{1+\tau^2\|M^{-1}\|^2}\,\varepsilon .
\]
Since \(e_0=0\), iterating gives
\[
    e_N
    \le
    \tau\sqrt{1+\tau^2\|M^{-1}\|^2}\,\varepsilon
    \sum_{j=0}^{N-1}(1+\tau L_\Phi)^j .
\]
Evaluating the geometric sum yields the result.
\end{proof}

\begin{corollary}[No hallucinated lateral escape]
\label{cor:no_hallucination}
Let \(P_\perp\) denote projection onto the lateral subspace. Suppose local
coordinates are scaffold-aligned so that longitudinal context forces do not generate
lateral motion to first order. If
\[
    \|P_\perp\nabla_qH_{\mathrm{ctx}}(q;z)\|\le \varepsilon_\perp
\]
on the rollout tube, then
\[
    \|P_\perp(q_N^{(2)}-q_N^{(1)})\|
    \le
    C_{\perp,N}\varepsilon_\perp,
\]
where
\[
    C_{\perp,N}
    =
    \tau^2
    \|P_\perp M^{-1}P_\perp^\top\|
    \frac{(1+\tau L_\Phi)^N-1}{\tau L_\Phi}.
\]
When \(\varepsilon_\perp=0\), the context channel cannot produce a lateral escape.
\end{corollary}

\begin{proof}
Project the context-induced position perturbation:
\[
    P_\perp \delta q_{t+1}
    =
    -\tau^2P_\perp M^{-1}\nabla_qH_{\mathrm{ctx}}(q_t;z_t).
\]
Under scaffold-aligned local coordinates, the first-order longitudinal--lateral cross
terms vanish, so
\[
    \|P_\perp \delta q_{t+1}\|
    \le
    \tau^2
    \|P_\perp M^{-1}P_\perp^\top\|
    \varepsilon_\perp .
\]
Propagating this one-step projected perturbation through the same Lipschitz recursion
as in Theorem~\ref{thm:scaffold_preservation} gives the claimed bound.
\end{proof}

\subsection{Selective risk deflection and risk reduction}

\begin{theorem}[Selective risk deflection and local risk reduction]
\label{thm:risk_deflection}
Let \(P_\perp\) project onto a feasible lateral subspace. Suppose that at a state \(q\),
\[
    \|P_\perp\nabla_q\widetilde r(\mathbf c;z)\|\ge \Delta,
\]
and suppose the projected hard-hazard force does not cancel the soft-risk direction by
more than \(\chi\):
\[
    \left\|
    P_\perp
    \left[
    \lambda_h b_{\mathrm{sp}}'(\phi(\mathbf c;z))\nabla_q\phi(\mathbf c;z)
    \right]
    \right\|
    \le \chi.
\]
Then
\[
    \|P_\perp F_{\mathrm{ctx}}(q;z)\|
    \ge
    \lambda_s\Delta-\chi .
\]
Thus, whenever \(\lambda_s\Delta>\chi\), the context channel is necessarily active in
the lateral subspace.

Assume further that \(\widetilde r\) has \(L_r\)-Lipschitz gradient,
\(M^{-1}\succeq m_M I\), and
\[
    \tau^2\lambda_s\|M^{-1}\|\le \frac{1}{L_r}.
\]
Then, ignoring shared geometry and dissipation terms up to \(O(\tau^3)\), the
one-step context-enriched perturbation satisfies
\[
    \widetilde r(q_{t+1}^{(2)})
    -
    \widetilde r(q_{t+1}^{(1)})
    \le
    -\frac12\tau^2\lambda_s m_M \Delta^2
    +
    O(\tau^3).
\]
Over any active window \(\mathcal A\) on which the same margin condition holds,
\[
    \sum_{t\in\mathcal A}
    \left[
    \widetilde r(q_{t+1}^{(2)})
    -
    \widetilde r(q_{t+1}^{(1)})
    \right]
    \le
    -\frac12|\mathcal A|\tau^2\lambda_s m_M \Delta^2
    +
    O(|\mathcal A|\tau^3).
\]
\end{theorem}

\begin{proof}
The context force is
\[
    F_{\mathrm{ctx}}
    =
    -\lambda_s\nabla_q\widetilde r
    -
    \lambda_h b_{\mathrm{sp}}'(\phi)\nabla_q\phi .
\]
Projecting onto the lateral subspace gives
\[
    P_\perp F_{\mathrm{ctx}}
    =
    -\lambda_sP_\perp\nabla_q\widetilde r
    -
    P_\perp\!\left[
    \lambda_h b_{\mathrm{sp}}'(\phi)\nabla_q\phi
    \right].
\]
By the reverse triangle inequality,
\[
    \|P_\perp F_{\mathrm{ctx}}\|
    \ge
    \lambda_s\|P_\perp\nabla_q\widetilde r\|
    -
    \left\|
    P_\perp\!\left[
    \lambda_h b_{\mathrm{sp}}'(\phi)\nabla_q\phi
    \right]
    \right\|.
\]
The assumptions imply
\[
    \|P_\perp F_{\mathrm{ctx}}\|\ge \lambda_s\Delta-\chi.
\]
This proves activation under \(\lambda_s\Delta>\chi\).

For the risk-reduction claim, isolate the soft-risk force
\[
    F_{\mathrm{soft}}=-\lambda_s\nabla_q\widetilde r(q).
\]
The position perturbation induced by this force in one semi-implicit step is
\[
    \delta q
    =
    -\tau^2\lambda_sM^{-1}\nabla_q\widetilde r(q).
\]
By \(L_r\)-smoothness,
\[
    \widetilde r(q+\delta q)
    \le
    \widetilde r(q)
    +
    \langle\nabla_q\widetilde r(q),\delta q\rangle
    +
    \frac{L_r}{2}\|\delta q\|^2 .
\]
The first-order term is
\[
    \langle\nabla_q\widetilde r(q),\delta q\rangle
    =
    -\tau^2\lambda_s
    \nabla_q\widetilde r(q)^\top
    M^{-1}
    \nabla_q\widetilde r(q).
\]
Since \(M^{-1}\succeq m_MI\) and
\(\|\nabla_q\widetilde r(q)\|\ge \Delta\) on the active projected direction,
\[
    \nabla_q\widetilde r(q)^\top M^{-1}\nabla_q\widetilde r(q)
    \ge
    m_M\Delta^2.
\]
Therefore
\[
    \langle\nabla_q\widetilde r(q),\delta q\rangle
    \le
    -\tau^2\lambda_s m_M\Delta^2.
\]
The second-order term is
\[
    \frac{L_r}{2}\|\delta q\|^2
    =
    \frac{L_r}{2}
    \tau^4\lambda_s^2
    \|M^{-1}\nabla_q\widetilde r(q)\|^2 .
\]
The step-size condition
\[
    \tau^2\lambda_s\|M^{-1}\|\le \frac{1}{L_r}
\]
ensures that this second-order term is controlled by the first-order decrease, yielding
\[
    \widetilde r(q+\delta q)-\widetilde r(q)
    \le
    -\frac12\tau^2\lambda_s m_M\Delta^2.
\]
The shared geometry and dissipative terms are identical in the geometry scaffold and context-enriched field
local comparison; by smoothness, their contribution to the relative risk difference is
\(O(\tau^3)\). Summing the one-step inequality over \(\mathcal A\) proves the final
claim.
\end{proof}

\subsection{CVaR gradient consistency}

The episode cost used in the main text is
\[
    J(\theta)
    =
    w_g\|q_T-q_g\|^2
    +
    w_\ell\sum_t\|q_{t+1}-q_t\|
    +
    w_r\sum_t \widetilde r(q_t)\|q_{t+1}-q_t\|
    +
    w_h\sum_t \mathbf 1\{\phi(q_t)<\epsilon\}.
\]
The training objective uses the Rockafellar--Uryasev representation of CVaR:
\[
    \operatorname{CVaR}_\alpha(J_\theta)
    =
    \min_{\eta\in\mathbb R}
    \left[
    \eta+
    \frac{1}{1-\alpha}\mathbb E(J_\theta-\eta)_+
    \right].
\]

\begin{theorem}[Detached empirical CVaR gives a consistent tail-gradient estimator]
\label{thm:cvar_consistency}
Let \(J_\theta(\xi)\) be the rollout cost induced by parameter \(\theta\) and rollout
randomness \(\xi\). Fix \(\alpha\in(0,1)\). Assume:
\begin{enumerate}
    \item \(J_\theta(\xi)\) is differentiable in \(\theta\) almost surely;
    \item \(\|\nabla_\theta J_\theta(\xi)\|\le G\) almost surely;
    \item the distribution of \(J_\theta\) has a continuous density near its
    \(\alpha\)-quantile;
    \item \(\mathbb P[J_\theta=\eta_\alpha(\theta)]=0\), where
    \(\eta_\alpha(\theta)=\operatorname{VaR}_\alpha(J_\theta)\).
\end{enumerate}
Then
\[
    \nabla_\theta \operatorname{CVaR}_\alpha(J_\theta)
    =
    \frac{1}{1-\alpha}
    \mathbb E
    \left[
    \nabla_\theta J_\theta(\xi)\,
    \mathbf 1\{J_\theta(\xi)\ge \eta_\alpha(\theta)\}
    \right].
\]
Given \(B\) i.i.d. rollouts and the detached empirical quantile
\[
    \widehat \eta
    =
    \widehat Q_\alpha(\{J_\theta^{(i)}\}_{i=1}^B),
\]
define
\[
    \widehat g_B(\theta)
    =
    \frac{1}{(1-\alpha)B}
    \sum_{i=1}^B
    \nabla_\theta J_\theta^{(i)}
    \mathbf 1\{J_\theta^{(i)}\ge \widehat \eta\}.
\]
Then
\[
    \widehat g_B(\theta)
    \to
    \nabla_\theta \operatorname{CVaR}_\alpha(J_\theta)
\]
almost surely as \(B\to\infty\), and
\[
    \|\widehat g_B(\theta)
    -
    \nabla_\theta \operatorname{CVaR}_\alpha(J_\theta)\|
    =
    O_p(B^{-1/2}).
\]
If \(J_\theta\in[0,J_{\max}]\), then the empirical CVaR value obeys
\[
    |\widehat{\operatorname{CVaR}}_\alpha-\operatorname{CVaR}_\alpha|
    =
    O\!\left(
    \frac{J_{\max}}{1-\alpha}
    \sqrt{\frac{\log(1/\delta)}{B}}
    \right)
\]
with probability at least \(1-\delta\).
\end{theorem}

\begin{proof}
The Rockafellar--Uryasev variational form gives
\[
    \operatorname{CVaR}_\alpha(J_\theta)
    =
    \min_{\eta\in\mathbb R}
    \left[
    \eta+
    \frac{1}{1-\alpha}\mathbb E(J_\theta-\eta)_+
    \right].
\]
At the minimizer,
\[
    \eta^\star=\eta_\alpha(\theta).
\]
By the envelope theorem, when differentiating the optimized value with respect to
\(\theta\), the derivative of the optimizer \(\eta^\star(\theta)\) does not appear. Hence
\[
    \nabla_\theta \operatorname{CVaR}_\alpha(J_\theta)
    =
    \frac{1}{1-\alpha}
    \nabla_\theta
    \mathbb E[(J_\theta-\eta^\star)_+].
\]
Since \(J_\theta\) is differentiable almost surely and the quantile boundary has
probability zero,
\[
    \nabla_\theta(J_\theta-\eta^\star)_+
    =
    \nabla_\theta J_\theta\,
    \mathbf 1\{J_\theta\ge \eta^\star\}
\]
almost surely. The bounded-gradient assumption permits interchanging derivative and
expectation, proving
\[
    \nabla_\theta \operatorname{CVaR}_\alpha(J_\theta)
    =
    \frac{1}{1-\alpha}
    \mathbb E
    \left[
    \nabla_\theta J_\theta
    \mathbf 1\{J_\theta\ge \eta_\alpha\}
    \right].
\]

For the empirical estimator, continuity of the distribution near \(\eta_\alpha\) implies
strong consistency of the empirical quantile:
\[
    \widehat \eta \to \eta_\alpha
\]
almost surely. Since there is no atom at the quantile,
\[
    \mathbf 1\{J_\theta^{(i)}\ge\widehat\eta\}
    \to
    \mathbf 1\{J_\theta^{(i)}\ge\eta_\alpha\}
\]
almost surely. The summands are bounded by \(G/(1-\alpha)\), so dominated
convergence and the strong law of large numbers give
\[
    \widehat g_B(\theta)
    \to
    \frac{1}{1-\alpha}
    \mathbb E
    \left[
    \nabla_\theta J_\theta
    \mathbf 1\{J_\theta\ge\eta_\alpha\}
    \right]
    =
    \nabla_\theta \operatorname{CVaR}_\alpha(J_\theta)
\]
almost surely.

For the rate, decompose
\[
    \widehat g_B-g
    =
    A_B+R_B,
\]
where
\[
    A_B
    =
    \frac{1}{(1-\alpha)B}
    \sum_{i=1}^B
    \left[
    \nabla J_i\mathbf 1\{J_i\ge\eta_\alpha\}
    -
    \mathbb E(\nabla J\mathbf 1\{J\ge\eta_\alpha\})
    \right]
\]
and
\[
    R_B
    =
    \frac{1}{(1-\alpha)B}
    \sum_{i=1}^B
    \nabla J_i
    \left[
    \mathbf 1\{J_i\ge\widehat\eta\}
    -
    \mathbf 1\{J_i\ge\eta_\alpha\}
    \right].
\]
The first term is a bounded empirical average, so \(A_B=O_p(B^{-1/2})\). The
second term is nonzero only when \(J_i\) lies between \(\widehat\eta\) and
\(\eta_\alpha\). The probability mass of this interval is
\(O_p(|\widehat\eta-\eta_\alpha|)\), and standard empirical quantile theory gives
\(|\widehat\eta-\eta_\alpha|=O_p(B^{-1/2})\). Hence \(R_B=O_p(B^{-1/2})\), proving
the gradient rate.

For bounded \(J_\theta\), the empirical Rockafellar--Uryasev objective is a bounded
empirical process indexed by \(\eta\). Uniform concentration over
\(\eta\in[0,J_{\max}]\), followed by comparison at the empirical and population
minimizers, yields the displayed
\[
    O\!\left(
    \frac{J_{\max}}{1-\alpha}
    \sqrt{\frac{\log(1/\delta)}{B}}
    \right)
\]
concentration rate.
\end{proof}

\subsection{Safety caveat: finite barriers are not CBF certificates}
\label{app:safety_caveat}

The hard-hazard term \(b_{\mathrm{sp}}(\phi)\) is a differentiable repulsive barrier. It
should not be stated as a formal forward-invariance certificate unless an additional
condition is imposed. In continuous time, forward invariance follows if the barrier
diverges at the unsafe boundary and total energy is bounded. In discrete time, one
needs either a sufficiently small step size that prevents crossing between samples, or
an explicit control-barrier-function projection/filter satisfying a discrete or
continuous CBF condition~\citep{Ames2017CBF}. Therefore, our theoretical claims are
scaffold preservation, dissipativity, selectivity, and local risk reduction under stated
regularity assumptions; they are not unconditional collision-avoidance guarantees.
